\documentclass[11pt,letterpaper]{article}

% ----------------------------------------------------------------------
% Layout
% ----------------------------------------------------------------------
\usepackage[margin=1in]{geometry}
\usepackage{setspace}
\usepackage{microtype}
\usepackage{enumitem}
\usepackage{caption}

% ----------------------------------------------------------------------
% Math and units
% ----------------------------------------------------------------------
\usepackage{amsmath,amssymb,amsthm}
\usepackage{siunitx}

% ----------------------------------------------------------------------
% Tables
% ----------------------------------------------------------------------
\usepackage{booktabs}

% ----------------------------------------------------------------------
% Code listings
% ----------------------------------------------------------------------
\usepackage{listings}
\usepackage{xcolor}

% ----------------------------------------------------------------------
% Links and bibliography
% ----------------------------------------------------------------------
\usepackage{hyperref}
\usepackage[numbers,sort&compress]{natbib}

% ----------------------------------------------------------------------
% Package configuration
% ----------------------------------------------------------------------
\sisetup{per-mode=symbol,detect-all}
\DeclareSIUnit\rpm{rpm}

\hypersetup{
  colorlinks = true,
  linkcolor  = black,
  citecolor  = blue!50!black,
  urlcolor   = blue!50!black,
}

\lstdefinestyle{code}{
  basicstyle       = \ttfamily\footnotesize,
  breaklines       = true,
  columns          = fullflexible,
  frame            = single,
  rulecolor        = \color{black!20},
  backgroundcolor  = \color{black!3},
  showstringspaces = false,
  commentstyle     = \itshape\color{black!55},
  keywordstyle     = \bfseries,
  morecomment      = [l]{//},
  morecomment      = [s]{/*}{*/},
}
\lstset{style=code}

% ----------------------------------------------------------------------
% Dual audience labels
% ----------------------------------------------------------------------
\newcommand{\eng}[1]{\paragraph{Physics.} #1}
\newcommand{\codepath}[1]{\paragraph{Software.} #1}

% ----------------------------------------------------------------------
% Title
% ----------------------------------------------------------------------
\title{%
  The UBCO 2026 EV Vehicle Model:\\
  Six-Degree-of-Freedom Dynamics, Quasi-Steady Laps,\\
  and Parameter Provenance%
}
\author{%
  UBCO Motorsports --- Vehicle Dynamics and Software\\
  Technical report: \texttt{Vehicle-Model-V2} \quad vehicle type \texttt{fsae-6dof}%
}
\date{August 2026}

\begin{document}

\maketitle

\begin{abstract}
This report specifies the physics and the software of the 2026 electric
Formula~SAE vehicle model used by UBCO Motorsports (OK Motorsports). It
is written for two audiences. \textbf{Vehicle dynamics, aero, powertrain,
and tire engineers} receive the Newton--Euler six-degree-of-freedom
(6DOF) chassis, tire forces, electric torque--speed envelope, aerodynamic
maps, G-G diagrams, and the three-pass quasi-steady-state (QSS) lap, with
SI units justified at each step. \textbf{Software engineers} receive the
matching objects in \texttt{fsae-6dof}, the OpenVEHICLE workbook, and the
Python QSS and heads-up display (HUD) stack, including a strict
separation between ordinary-differential-equation (ODE) states and
reconstructed dashboard channels.

The C++ plant is a Formula~SAE specialisation of Manzanero's
\texttt{fastest-lap} Newton--Euler engine~\cite{manzanero-fastestlap}.
The parameter books and the QSS racing-line method continue the club
Vehicle Dynamics (VD) MATLAB OpenLAP / OpenVEHICLE workflow developed
under Brenden Jones~\cite{jones-linkedin,openlap,lopesbandy-linkedin}.
If a number is edited in Studio, this document states which equation it
enters and why.
\end{abstract}

\tableofcontents
\newpage

% ========================================================================
\section{How to read this document}
% ========================================================================

\paragraph{Physics} paragraphs state the model, assumptions, and units.
\paragraph{Software} paragraphs name the repository files. They are
pointers into \texttt{Vehicle-Model-V2}, not a second implementation.

New software members: start with Section~\ref{sec:lineage},
Section~\ref{sec:two-models}, Section~\ref{sec:code}, and
Appendix~\ref{app:files}. Dynamics leads: treat the numbered equations
as the specification and read Section~\ref{sec:limits} before using
outputs in a design review.

Two pipelines share one car file:

\begin{enumerate}[leftmargin=1.6em]
  \item \textbf{Transient 6DOF} (C++): Newton--Euler sprung body, four
        Pacejka tires, electric-vehicle (EV) envelope. Used for
        numerical G-G diagrams (\texttt{gg\_diagram}) and, later,
        optimal lap time.
  \item \textbf{QSS lap and HUD} (Python): OpenTRACK mesh, a G-G table,
        three kinematic passes, then bicycle-model post-processing for
        steer, sideslip, and lumped tire loads. This is the Studio path
        on Cloudflare (Pyodide) and locally.
\end{enumerate}

They are not the same fidelity. Treating a Studio lap-time delta as a
6DOF G-G delta is a modelling error. Section~\ref{sec:two-models}
states the distinction formally.

% ========================================================================
\section{Previous 2DOF MATLAB model and OpenLAP}
\label{sec:lineage}
% ========================================================================

UBCO Motorsports is the University of British Columbia Okanagan
Formula~SAE team~\cite{okmotorsports}. Vehicle Dynamics owns how the car
uses the road: mass distribution, springs, anti-roll bars, kinematics,
tires, aerodynamic balance, and the simulations that convert those
numbers into predicted skidpad, acceleration, and endurance performance.

\subsection{Technical work under Brenden Jones}
\label{sec:jones}

From 2023, Brenden Jones has been Vehicle Dynamics and Software Lead
for OK Motorsports~\cite{jones-linkedin}. The technical product of that
role, continued from earlier OpenLAP use on the subteam
~\cite{lopesbandy-linkedin}, is a MATLAB toolchain that:

\begin{itemize}[leftmargin=1.6em]
  \item stored one official vehicle book (OpenVEHICLE) and one official
        track book (OpenTRACK);
  \item built a speed-dependent G-G (GGV) envelope from mass, power,
        grip, and aero;
  \item ran a quasi-steady lap along the racing line to predict lap
        time and longitudinal/lateral acceleration histories;
  \item reconstructed steer, body sideslip, and kinematic tire loads
        with a planar two-degree-of-freedom (2DOF) bicycle model;
  \item kept a separate 2DOF quarter-car for ride (sprung and unsprung
        vertical motion), which is still the ride tool.
\end{itemize}

Vehicle-Model-V2 does not replace that parameter ownership. It reads
the same workbooks and keeps QSS as the operational lap method, while
the plant used for G-G (and later optimal lap time) is upgraded from
planar 2DOF to sprung 6DOF.

\subsection{MATLAB architecture: OpenVEHICLE, OpenTRACK, OpenLAP}
\label{sec:matlab-arch}

OpenLAP is a three-program MATLAB stack~\cite{openlap}:

\begin{center}
\begin{tabular}{@{}llp{7.4cm}@{}}
\toprule
Program & Role & Inputs / outputs \\
\midrule
OpenVEHICLE & Car database & Sheets \texttt{Info}, \texttt{Torque Curve},
  \texttt{FastestLap}. Columns
  \texttt{Category $|$ Description $|$ Value $|$ Unit $|$ Comment}. \\
OpenTRACK & Track database & \texttt{Info} plus \texttt{Shape}:
  segment type (\texttt{Straight} / \texttt{Left} / \texttt{Right}),
  length in metres, corner radius in metres. The polyline \emph{is}
  the racing line; there is no corridor width. \\
OpenLAP & Lap solver & GGV table, three-pass $v(s)$, then channel
  reconstruction (throttle, brake, steer, $\beta$, axle loads). \\
\bottomrule
\end{tabular}
\end{center}

\codepath{\texttt{examples/python/fsae/openvehicle\_xlsx.py},
\texttt{opentrack\_xlsx.py}, \texttt{xlsx\_kit.py}. Studio patches
\texttt{Info} / \texttt{FastestLap} by Description, then writes
\texttt{ubco-2026-ev.xml} for C++.}

\subsection{Planar 2DOF bicycle (yaw and sideslip)}
\label{sec:2dof-bicycle}

The MATLAB \emph{handling} model is a single-track bicycle
~\cite{gillespie,milliken}: two planar degrees of freedom after speed
is treated as a prescribed (or slowly varying) parameter. With
longitudinal speed $u$ (\si{\metre\per\second}), lateral speed $v$,
yaw rate $r$ (\si{\radian\per\second}), wheelbase $L=a+|b|$, and linear
cornering stiffnesses $C_F,C_R$ (\si{\newton\per\radian}),
\begin{align}
  m(\dot v + u r) &= F_{yF} + F_{yR}, \label{eq:2dof-lat}\\
  I_{zz}\dot r &= a F_{yF} + b F_{yR}, \label{eq:2dof-yaw}
\end{align}
where $b$ is negative if measured rearward from the centre of gravity
(CG). Small-angle slip:
\begin{equation}
  \alpha_F = \delta - \frac{v + a r}{u},\qquad
  \alpha_R = -\frac{v + b r}{u},\qquad
  F_{yF}=C_F\alpha_F,\quad F_{yR}=C_R\alpha_R.
\end{equation}
In OpenLAP, this bicycle is used primarily as a \emph{post-process}:
path curvature $\kappa(s)$ and speed $v(s)$ are already known from QSS,
so centripetal demand $m v^2\kappa$ is inverted for extra steer and
body sideslip $\beta=v/u$ (Section~\ref{sec:hud}). The bicycle is not
integrated as a yaw-inertia ODE along the track.

\eng{Equations~\eqref{eq:2dof-lat}--\eqref{eq:2dof-yaw} give
understeer gradient and a steer trace that MATLAB users already trust.
They cannot produce left/right $F_z$ from springs, roll-camber, or
aero that depends on ride height.}

\subsection{Quasi-steady lap in MATLAB}
\label{sec:matlab-qss}

OpenLAP does not integrate the bicycle in time for endurance lap time.
It uses the same three-pass construction now in Python
(Section~\ref{sec:qss}):

\begin{enumerate}[leftmargin=1.6em]
  \item Build or import a GGV table: at each speed, the admissible
        $(a_x,a_y)$ set, with $a_x,a_y$ stored in units of $g$.
  \item Corner limit: $v_{\max}(s)=\sqrt{a_y^{\max}(v)/|\kappa(s)|}$.
  \item Forward pass: accelerate no harder than $a_x^{\max}(v,a_y)$.
  \item Backward pass: brake no harder than $|a_x^{\min}|$.
  \item $v=\min(v_{\max},v_{\mathrm{acc}},v_{\mathrm{brk}})$, then
        $t=\int \mathrm{d}s/v$.
\end{enumerate}

Kinematic tire loads after the lap use lumped transfer
$\Delta F_z \propto m a h / t$ (Section~\ref{sec:hud}), not a sprung
roll ODE.

\subsection{Quarter-car 2DOF (ride)}
\label{sec:quarter}

Ride was, and remains, a separate MATLAB 2DOF quarter-car: sprung mass
$m_s$ and unsprung mass $m_u$,
\begin{align}
  m_s\ddot z_s &= -k_s(z_s-z_u)-c_s(\dot z_s-\dot z_u),\\
  m_u\ddot z_u &= k_s(z_s-z_u)+c_s(\dot z_s-\dot z_u)
                 -k_t(z_u-z_r),
\end{align}
with road $z_r(t)$ and tire stiffness $k_t$. That model is the correct
tool for wheel hop and damper curves. The 6DOF plant in this report
does \emph{not} absorb it: corners are algebraic (Section~\ref{sec:corner}).

\subsection{Why 6DOF is the upgrade}
\label{sec:upgrade}

\begin{center}
\begin{tabular}{@{}p{0.30\textwidth}p{0.32\textwidth}p{0.32\textwidth}@{}}
\toprule
Capability & MATLAB 2DOF / OpenLAP & \texttt{fsae-6dof} \\
\midrule
Rigid-body motion
  & Planar yaw + sideslip; ride in a separate quarter-car
  & Coupled surge, sway, heave, roll, pitch, yaw. \\
Normal load
  & Algebraic $m a h/t$ after QSS
  & $F_z$ from series springs, ARB, and Euler roll/pitch. \\
Aero
  & Speed in the GGV table; little or no $z,\mu$ coupling
  & $C_L(z,\mu)$, $C_D(z,\mu)$ at a pressure centre. \\
Tires
  & $\mu$ or linear $C_F,C_R$
  & Combined-slip Pacejka; $\kappa$ a state per wheel. \\
Powertrain
  & Torque curve / power in OpenVEHICLE
  & EV envelope, regen, viscous LSD, independent rear spins. \\
Lap time (Studio)
  & OpenLAP three-pass
  & Same three-pass, optional 6DOF G-G as the table. \\
Ride / wheel hop
  & Quarter-car 2DOF
  & Unchanged; still MATLAB. \\
\bottomrule
\end{tabular}
\end{center}

The upgrade is coupling. A Formula~SAE car generates lateral
acceleration from $F_z$ and camber; $F_z$ comes from roll, heave,
pitch, and downforce; downforce depends on ride height and pitch. A
planar bicycle with post-hoc load transfer cannot close that loop.
The 6DOF plant can, which is why G-G and optimal lap time use
\texttt{fsae-6dof} rather than the MATLAB ODE.

Continuity for the club: OpenVEHICLE / OpenTRACK files remain the
parameter source, QSS remains the daily lap method, and HUD channels
remain OpenLAP-shaped so existing MATLAB plots are comparable.

\eng{Trust the MATLAB car file as the official numbers. This stack
reads the same cells; it does not maintain a second mass or power.}

% ========================================================================
\section{What the model borrows (and what it does not)}
\label{sec:borrows}
% ========================================================================

\begin{table}[ht]
\centering
\caption{Intellectual inventory.}
\begin{tabular}{@{}llp{7.2cm}@{}}
\toprule
Layer & Source & What we took \\
\midrule
Newton--Euler 6DOF chassis
  & fastest-lap~\cite{manzanero-fastestlap}
  & $\boldsymbol{\omega}\times\mathbf{v}$ in the body frame; Euler
    equations with small roll/pitch; kart-style algebraic corners. \\
Pacejka tires
  & Pacejka~\cite{pacejka}, fastest-lap \texttt{Tire\_pacejka\_std}
  & Combined slip $F_x,F_y$ vs $\kappa,\alpha,F_z$. Shape coefficients
    remain placeholders pending TTC~\cite{ttc}. \\
G-G thinking
  & Milliken~\cite{milliken}
  & Admissible $(a_x,a_y)$ at a given speed; friction ellipse. \\
Bicycle steer post-process
  & Gillespie~\cite{gillespie}, OpenLAP~\cite{openlap}
  & HUD $\delta,\beta$ from $v,\kappa,C_F,C_R$ --- not 6DOF states. \\
QSS three-pass lap
  & OpenLAP~\cite{openlap}
  & $v_{\max}$, forward acceleration, backward braking along the
    racing line. \\
EV envelope
  & FSAE rules~\cite{fsae-rules}
  & $T=\min(T_{\mathrm{peak}}, P_{\max}/\omega)$ after the gearbox. \\
Aero maps
  & Katz-style linearisation~\cite{katz}
  & $C_L(z,\mu)$, $C_D(z,\mu)$ about a pressure centre. \\
Parameters
  & VD OpenVEHICLE workflow~\cite{jones-linkedin,openlap}
  & Mass, tracks, wheel rates, $C_L/C_D$, 80\,\si{\kilo\watt}, 2019
    endurance map. \\
\bottomrule
\end{tabular}
\end{table}

The MATLAB OpenLAP ODE was not ported into C++. Unsprung masses, 3D
wishbone inverse kinematics, battery state of charge, ABS, and track
banking are out of scope (design specs
\texttt{2026-08-14-fsae-6dof-design.md} and
\texttt{2026-08-14-fsae-further-ideas-design.md}).

% ========================================================================
\section{Two models, one car file}
\label{sec:two-models}
% ========================================================================

\begin{center}
\begin{tabular}{@{}p{0.47\textwidth}p{0.47\textwidth}@{}}
\toprule
\textbf{6DOF plant (G-G / optimal lap time)} &
\textbf{QSS + HUD (Studio)} \\
\midrule
States: body $u,v,r,z,\phi,\mu$ and rates;
four wheel spins; four tire temperatures.
Forces from Pacejka at each contact patch.
Aero at the pressure centre.
EV torque on the rear axle.
&
States along the track: $s$, $v(s)$ only.
$a_y = v^2 \kappa / g$, $a_x$ from
$v^2$ differences.
Steer, $\beta$, $F_z$, TPS/BPS are
reconstructed after the lap. \\
\midrule
File: \texttt{src/core/vehicles/fsae6dof.h} &
File: \texttt{examples/python/fsae/qss\_lap.py} \\
\bottomrule
\end{tabular}
\end{center}

\eng{A ride-rate change in the 6DOF plant can move G-G through $F_z$
and camber. The same change in Studio QSS with the vehicle-parameter
envelope affects lap time only through mass, $\mu$, power, and aero,
not through Pacejka. Those two deltas are different tests.}

% ========================================================================
\section{Why the system is six degree-of-freedom}
\label{sec:why6}
% ========================================================================

A free rigid body in three-dimensional space has six degrees of
freedom: three translations and three rotations. For the sprung mass
the implementation uses a road-aligned chassis
frame~\cite{manzanero-fastestlap,genta}:

\begin{center}
\begin{tabular}{@{}lll@{}}
\toprule
DOF & Symbol & Unit \\
\midrule
Surge (longitudinal) & $u$ or $v_x$          & \si{\metre\per\second} \\
Sway (lateral)        & $v$ or $v_y$          & \si{\metre\per\second} \\
Heave (vertical)      & $z$                   & \si{\metre} \\
Roll                   & $\phi$                & \si{\radian} \\
Pitch                  & $\mu$                 & \si{\radian} \\
Yaw & $\psi$ (road) / $r=\dot\psi$ (body yaw rate)
    & \si{\radian}, \si{\radian\per\second} \\
\bottomrule
\end{tabular}
\end{center}

\paragraph{Why the plant is not a bicycle.}
A planar bicycle (yaw and lateral velocity only) cannot provide:
(i)~normal load from heave, roll, pitch, and the series spring
network; (ii)~aerodynamics that depend on ride height $z$ and pitch
$\mu$; (iii)~left/right camber from roll; (iv)~independent rear wheel
speeds for a viscous limited-slip differential (LSD). Formula~SAE
performance is dominated by load transfer and aero, which is why the
plant is 6DOF.

\paragraph{Why not 14DOF with unsprung masses.}
Each unsprung corner adds vertical (and often spin) states. The MATLAB
2DOF quarter-car remains the ride tool (Section~\ref{sec:quarter}).
The 6DOF plant treats the corner as an algebraic spring--damper into
the tire, with hub motion slaved to $z,\phi,\mu$: fewer states for G-G
and optimal-lap nonlinear programs (NLPs), at the cost of wheel hop.

\paragraph{Internal states that are not the six rigid DOF.}
Four dimensionless longitudinal slips $\kappa$ (or equivalently wheel
spins) and four carcass temperatures $T_i$ are tire and thermal
dynamics attached to the same body.

\codepath{Chassis states in \texttt{chassis\_car\_6dof.h}: $z,\phi,\mu$
and $\dot z,\dot\phi,\dot\mu$. The planar chassis base supplies $u,v,r$.
The FSAE chassis adds $T_{\mathrm{FL,FR,RL,RR}}$. Each FSAE axle adds
left/right $\kappa$ and spin momentum.}

% ========================================================================
\section{Units, constants, and sign conventions}
\label{sec:units}
% ========================================================================

Unless a table says otherwise, the plant is SI.

\begin{table}[ht]
\centering
\caption{Constants and unit justification.}
\begin{tabular}{@{}llll@{}}
\toprule
Symbol      & Value              & SI unit & Justification \\
\midrule
$g$         & $9.81$             & \si{\metre\per\second\squared}
            & Standard gravity. HUD $a_x,a_y$ in ``$g$'' are $a/g$
              (dimensionless). \\
$m$         & $277.2$            & \si{\kilo\gram}
            & Inertia in $F=ma$. Includes driver in the VD book. \\
$\rho$      & $1.2$              & \si{\kilo\gram\per\metre\cubed}
            & Near-ISA sea-level density in
              $D=\tfrac12\rho C_D A v^2$. \\
$A$         & $1.14$             & \si{\metre\squared}
            & Frontal area for aerodynamic force. \\
$P_{\max}$  & $80\times10^3$     & \si{\watt}
            & FSAE EV cap~\cite{fsae-rules}. Stored as \si{\kilo\watt}
              in the workbook. \\
$T_{\mathrm{peak}}$ & $240$      & \si{\newton\metre}
            & Motor torque; $\mathrm{N\cdot m} = \mathrm{J}$. \\
$k_w$       & $1.3\times10^5$    & \si{\newton\per\metre}
            & Wheel rate: force per metre of wheel travel. \\
$c_s$       & $1500$             & \si{\newton\second\per\metre}
            & Damper: force per \si{\metre\per\second}. \\
$I_{zz}$    & $110$              & \si{\kilo\gram\metre\squared}
            & Yaw inertia; $M = I\alpha$. \\
\bottomrule
\end{tabular}
\end{table}

\paragraph{Workbook millimetres.}
OpenVEHICLE stores wheelbase, tracks, CG height, and tire radius in
\si{\milli\metre} because that is the MATLAB convention. The importer
divides by $1000$ before XML. Example: wheelbase
\SI{1620}{\milli\metre} $= \SI{1.620}{\metre}$. CG in XML is
\texttt{com.z = -h} with $h=\SI{0.25}{\metre}$ so geometric height is
positive up in the fastest-lap car frame.

\paragraph{OpenVEHICLE $C_L$ / $C_D$ signs.}
MATLAB OpenVEHICLE uses $C_L>0$ as lift and, in the club book,
typically negative $C_D$. fastest-lap uses $+C_L$ as downforce and
$+C_D$ as drag. The importer flips sign:
\[
  C_{L,\mathrm{FL}} = -C_{L,\mathrm{OV}}\cdot f_{C_L},\qquad
  C_{D,\mathrm{FL}} = -C_{D,\mathrm{OV}}\cdot f_{C_D}.
\]
Club defaults $C_{L,\mathrm{OV}}=-3.77$, $C_{D,\mathrm{OV}}=-1.4$ become
downforce $3.77$ and drag $1.4$ in XML.

\paragraph{Accelerations in the HUD.}
QSS stores $a_x,a_y$ in $g$. Multiply by
$9.81\,\si{\metre\per\second\squared}$ for SI. Speed in plots is
$v\cdot 3.6$ for \si{\kilo\metre\per\hour}.

% ========================================================================
\section{Newton--Euler 6DOF: governing equations}
\label{sec:newton}
% ========================================================================

In an inertial frame, $\sum\mathbf{F}=m\mathbf{a}$ and
$\sum\mathbf{M}=\dot{\mathbf{H}}$. The chassis is written in a
\textbf{road-projected body frame} that yaws with the car. Roll and
pitch are small extra rotations that move the axle frames; they are
not full direction-cosine updates~\cite{manzanero-fastestlap}. The
chassis frame remains parallel to the road projection of the CG frame
(\texttt{chassis\_car\_6dof.h}).

\paragraph{Linear momentum (body frame).}
For velocity $\mathbf{v}$ measured in a rotating frame with angular
velocity $\boldsymbol{\omega}$,
\begin{equation}
  \dot{\mathbf{v}}_{\mathrm{inertial}}
  = \dot{\mathbf{v}}_{\mathrm{body}} + \boldsymbol{\omega}\times\mathbf{v}.
  \label{eq:transport}
\end{equation}
Newton's law in the body frame is therefore
\begin{equation}
  \dot{\mathbf{v}}_{\mathrm{body}}
  = -\boldsymbol{\omega}\times\mathbf{v} + \frac{\mathbf{F}_{\mathrm{total}}}{m}.
  \label{eq:newton}
\end{equation}
The code names $\boldsymbol{\omega}\times\mathbf{v}$ as \texttt{Newton\_lhs}:

\begin{lstlisting}[language=C++,caption={Transport theorem, \texttt{chassis\_car\_6dof.hpp}.}]
return cross(omega, vel);  // Newton_lhs = omega x v
// dvdt = -Newton_lhs + F_total / m
\end{lstlisting}

$\mathbf{F}_{\mathrm{total}}$ is tire force (rotated into the road
frame) plus $mg$ on $Z$ plus scaled aero. Gravity is $+m g_0$ on $Z$ in
the fastest-lap road-frame convention.

\paragraph{Angular momentum.}
Euler's equation for a rigid body,
\begin{equation}
  \mathbf{I}\dot{\boldsymbol{\omega}}
  + \boldsymbol{\omega}\times(\mathbf{I}\boldsymbol{\omega})
  = \mathbf{M}_{\mathrm{total}},
  \label{eq:euler}
\end{equation}
is specialised for small roll $\phi$ and pitch $\mu$ with yaw rate
$\omega\equiv r$. Moments are taken about the road contact geometry,
not a free-space CG, which is why $m h$ coupling terms appear. The
code solves
\[
  \mathbf{E}_m\,\ddot{\boldsymbol{\varphi}}
  = -\mathbf{E}_{\mathrm{lhs}} + \mathbf{M}_{\mathrm{total}},
\]
where $\ddot{\boldsymbol{\varphi}}=(\ddot\phi,\ddot\mu,\dot r)$ and (with
$h=-\texttt{com.z}$, $z$ the heave state)
\begin{align}
  E_{\mathrm{lhs},x}
  &= m\bigl(\dot v\, h + (h-z)\,r u\bigr)
     - (I_{yy}-I_{zz}+I_{xx}) r \dot\mu
     - (I_{yy}-I_{zz}) r^2 \phi,\\
  E_{\mathrm{lhs},y}
  &= m\bigl((z-h)\dot u + h r v\bigr)
     + (I_{yy}-I_{zz}+I_{xx}) r \dot\phi
     - \bigl((I_{xx}-I_{zz})\mu + I_{xz}\bigr) r^2,\\
  E_{\mathrm{lhs},z}
  &= 2 I_{zx} \dot\mu\, r,
\end{align}
\begin{equation}
  \mathbf{E}_m =
  \begin{pmatrix}
    I_{xx} & 0      & -(I_{xx}-I_{zz})\mu - I_{xz} \\
    0      & I_{yy} & (I_{yy}-I_{zz})\phi \\
    -I_{xz} & 0     & I_{zz} + 2 I_{xz} \mu
  \end{pmatrix}.
\end{equation}
Those entries follow from a small-angle linearisation that mixes
roll and pitch into the yaw inertia row, keeping the chassis frame
parallel to the road.

Moments come from axle forces at
$\mathbf{x}_{\mathrm{front}},\mathbf{x}_{\mathrm{rear}}$ and aero at the
pressure centre plus current heave:
\begin{equation}
  \mathbf{M}
  = \mathbf{T}_f + \mathbf{x}_f\times\mathbf{F}_f
  + \mathbf{T}_r + \mathbf{x}_r\times\mathbf{F}_r
  + \bigl(\mathbf{x}_{\mathrm{pc}} + z\,\mathbf{e}_z\bigr)
    \times\mathbf{F}_{\mathrm{aero}}.
  \label{eq:moment}
\end{equation}

\eng{Roll moment from lateral force is $F_y h$ plus
$\mathbf{x}\times\mathbf{F}$ of left/right $F_z$. Load transfer is a
consequence of the dynamics, not an independent spreadsheet.}

\paragraph{Why G-G holds $\dot z=\dot\phi=\dot\mu=0$.}
A G-G point is a steady acceleration at fixed speed. The NLP asks for
controls and slips such that heave, roll, pitch, and yaw residuals
match a prescribed $(a_x,a_y)$. Tire temperatures are held at ambient
so thermal transients do not enter the envelope. See
\texttt{steady\_state\_equations} in \texttt{fsae6dof.h}.

% ========================================================================
\section{Corner compliance: why $F_z = k_t w + c_s \dot s$}
\label{sec:corner}
% ========================================================================

Each corner is not a separate unsprung mass. Vertical tire compression
$w$ is an algebraic function of heave, roll, and (front) steer jacking,
using the kart-style series network~\cite{manzanero-fastestlap}.

Let $k_t$ be tire vertical stiffness (\si{\newton\per\metre}), $k_w$
wheel rate, $k_{\mathrm{AR}}$ anti-roll stiffness referred to the axle.
Two reduction coefficients appear in the kart-style network. The code
names them \texttt{sym\_stiffness} and \texttt{assym\_stiffness}. The
first is a compliance; the second is dimensionless as used in $w_{L,R}$:
\begin{equation}
  C_{\mathrm{sym}} = \frac{1}{k_w + k_t}
  \qquad [\si{\metre\per\newton}],
\end{equation}
\begin{equation}
  C_{\mathrm{asym}}
  = \frac{2 k_{\mathrm{AR}} k_t C_{\mathrm{sym}}}{2 k_{\mathrm{AR}} + k_w + k_t}
  \qquad [-].
\end{equation}
With axle height $z_{\mathrm{axle}}$ (metres, road frame $Z$) and
unloaded radius $R_0$ (metres),
\[
  \delta_{\mathrm{sym}} = z_{\mathrm{axle}} + R_0,\qquad
  \delta_{\mathrm{asym}} = \tfrac12 t\,\phi
  \quad\text{(plus front jacking $\beta_R\delta$ if steering)},
\]
the algebraic tire compressions (metres) are
\begin{align}
  w_L &= k_w C_{\mathrm{sym}}(\delta_{\mathrm{sym}}-\delta_{\mathrm{asym}})
        - \delta_{\mathrm{asym}} C_{\mathrm{asym}},\\
  w_R &= k_w C_{\mathrm{sym}}(\delta_{\mathrm{sym}}+\delta_{\mathrm{asym}})
        + \delta_{\mathrm{asym}} C_{\mathrm{asym}}.
\end{align}
The factor $k_w C_{\mathrm{sym}} = k_w/(k_w+k_t)$ is the series-spring
share: only a fraction of chassis travel becomes tire squash. Slider
deflections $s_L,s_R$ reconstruct damper length; rates $\dot s$ follow
by replacing $\delta$ with $\dot\delta$. Contact load is
\begin{equation}
  F_{z,i} = \mathrm{smooth\_pos}\bigl(k_t w_i + c_s \dot s_i,\, 1\bigr),
  \label{eq:fz}
\end{equation}
so the damper is not an extra ODE. \texttt{smooth\_pos} enforces
$F_z\ge 0$. $k_t=\SI{240000}{\newton\per\metre}$ in the UBCO XML.

\paragraph{Units.}
$k_t w$: $(\mathrm{N/m})\cdot\mathrm{m} = \mathrm{N}$.
$c_s \dot s$: $(\mathrm{N\cdot s/m})\cdot(\mathrm{m/s}) = \mathrm{N}$.

\paragraph{UBCO numbers (VD book $\to$ XML).}
Front $k_w=\SI{130000}{\newton\per\metre}$,
$k_{\mathrm{AR}}=\SI{11260}{\newton\per\metre}$,
$c_s=\SI{1500}{\newton\second\per\metre}$; rear $143900$, $11160$,
$1400$. Tracks $t_f=\SI{1.225}{\metre}$, $t_r=\SI{1.220}{\metre}$.

\eng{A stiffer anti-roll bar increases $C_{\mathrm{asym}}$ and therefore
$F_z$ transfer at a given roll $\phi$. Roll itself is solved from
Euler; transfer cannot be prescribed independently of $h$, $t$, and
$k_{\mathrm{AR}}$.}

% ========================================================================
\section{Tires: Pacejka, $\kappa$, $\alpha$, and TTC status}
\label{sec:tires}
% ========================================================================

The contact patch does not obey Coulomb friction as a hard diamond.
Pacejka's Magic Formula~\cite{pacejka} gives smooth $F_x(\kappa,F_z)$
and $F_y(\alpha,F_z)$ and a combined-slip reduction. The sine form used
in \texttt{Tire\_pacejka\_std} is
\begin{equation}
  F = D\sin\bigl(C\arctan(B x - E(B x - \arctan B x))\bigr),
\end{equation}
with $x$ the slip ($\kappa$ or $\alpha$), $D=\mu F_z$ the peak
(\si{\newton}), $C$ the shape, $B$ the stiffness factor, $E$ the
curvature. Combined slip then scales the pure curves by cosine
weighting factors $r_{Bx},r_{Cx}$ (and lateral analogues), so peak
$F_x$ and peak $F_y$ cannot occur together.

Longitudinal slip (dimensionless)
\begin{equation}
  \kappa \approx \frac{\omega R - v_x}{v_x}
\end{equation}
is a state on each FSAE wheel (F1-3DOF style), not an algebraic
function of throttle. The update path is
\texttt{update\_from\_kappa}$(\kappa\cdot\kappa_{\max}(F_z))$, which
lets a G-G NLP place slips on the friction ellipse.

\paragraph{Relation of $p_{Dx1}$ and $p_{Dy1}$ to workbook $\mu$.}
XML \texttt{pDx1}$=0.9$ and \texttt{pDy1}$=1.5$ match OpenVEHICLE
$\mu_x,\mu_y$. They are peak friction coefficients (dimensionless):
$D_x \approx p_{Dx1} F_z$. They are not a fitted Pacejka shape until a
TTC identification replaces $B,C,E,K$.

\paragraph{Thermal grip.}
One node per tire~\cite{milliken}:
\begin{equation}
  C\frac{\mathrm{d}T}{\mathrm{d}t} = |\mathcal{D}| - h(T-T_{\mathrm{amb}}),
  \label{eq:thermal}
\end{equation}
with $C=\SI{900}{\joule\per\kelvin}$ (carcass thermal capacity),
$h=\SI{15}{\watt\per\kelvin}$, $T_{\mathrm{amb}}=\SI{298.15}{\kelvin}$,
$T_{\mathrm{opt}}=\SI{353.15}{\kelvin}$. Dissipation $\mathcal{D}$ is
Pacejka sliding power (\si{\watt}). Grip on tangential forces only:
\begin{equation}
  \mu_{\mathrm{scale}}
  = \mathrm{clip}\Bigl(1 - k_\mu \bigl((T-T_{\mathrm{opt}})/T_{\mathrm{opt}}\bigr)^2,\;
    0.5,\; 1.1\Bigr),
\end{equation}
$k_\mu=0.3$. G-G holds $T=T_{\mathrm{amb}}$ (isothermal).

\paragraph{Data status.}
Pacejka \emph{shape} coefficients are placeholders until a Hoosier TTC
fit~\cite{ttc}. Interim QSS envelopes use $\mu_x,\mu_y$ from the
OpenVEHICLE Info sheet (defaults $0.9$ and $1.5$). When TTC data are
fitted, replace the Pacejka block. Do not treat a Studio $\mu$ edit as
a tire model.

% ========================================================================
\section[Electric powertrain torque-speed envelope]%
        {Electric powertrain: why $T=\min(T_{\mathrm{peak}},P_{\max}/\omega)$}
\label{sec:ev}
% ========================================================================

Below base speed the inverter is torque-capped; above base speed it is
power-capped~\cite{fsae-rules}. With \texttt{peak-torque} present,
\texttt{Engine::operator()} evaluates
\begin{equation}
  \omega_m = n_g\,\omega_{\mathrm{axle}},\qquad
  T_{\mathrm{motor}}
  = u_{\mathrm{cmd}}\min\bigl(T_{\mathrm{peak}},\;
    P_{\max}/(\omega_m+\epsilon)\bigr),
  \label{eq:envelope}
\end{equation}
then returns axle torque $T_{\mathrm{axle}}=n_g T_{\mathrm{motor}}$. XML
stores $P_{\max}$ in \si{\kilo\watt}; C++ multiplies by $10^3$ to obtain
\si{\watt}. The identity $P=T\omega$ requires
$[\mathrm{W}]=[\mathrm{N\cdot m}][\mathrm{rad/s}]$.

Club numbers: $T_{\mathrm{peak}}=\SI{240}{\newton\metre}$,
$P_{\max}=\SI{80}{\kilo\watt}$, $n_g=4.8$. Base motor speed
\[
  \omega_{\mathrm{base}} = P_{\max}/T_{\mathrm{peak}}
  = 80000/240 \approx \SI{333}{\radian\per\second}
  \approx \SI{3180}{\rpm}.
\]
The argument \texttt{angular\_speed} is the mean of left and right
wheel spins (\si{\radian\per\second}), not motor rpm.

Throttle $u\in[-1,1]$:
\begin{align}
  u>0 &\Rightarrow T = u\cdot T_{\mathrm{env}}(\omega),\\
  u<0 &\Rightarrow T = u_{\mathrm{regen}}\cdot T_{\mathrm{env}},\quad
  u_{\mathrm{regen}} = \rho_{\mathrm{regen}} |u|,
\end{align}
with $\rho_{\mathrm{regen}}=0.3$, plus friction brakes
$T_{\mathrm{brake}} = T_{\mathrm{brake,max}}\cdot\mathrm{pos}(-u)\cdot b$
on each axle, $b$ = brake bias (1 = all front). Default bias $0.53$.

Rear viscous LSD:
\begin{equation}
  T_{L,R} = \tfrac12 T_{\mathrm{eng}} \mp k_d(\omega_L-\omega_R),
\end{equation}
$k_d=\SI{80}{\newton\metre\second\per\radian}$.

Wheel spin ODE (left):
\begin{equation}
  I_w \dot\omega_L = T_L + T_{\mathrm{tire},x,L},
\end{equation}
$I_w=\SI{0.5}{\kilo\gram\metre\squared}$ per axle inertia entry.

\codepath{\texttt{axle\_car\_6dof\_fsae.hpp} engine torque split;
\texttt{Engine} envelope when XML contains \texttt{peak-torque}.}

% ========================================================================
\section{Aerodynamics}
\label{sec:aero}
% ========================================================================

At the pressure centre $\mathbf{x}_{\mathrm{pc}}$,
\begin{equation}
  \mathbf{F}_{\mathrm{aero}}
  = \tfrac12\rho v^2 A \bigl( C_D\,\mathbf{e}_{\mathrm{drag}}
  + C_L\,\mathbf{e}_{\mathrm{lift}}\bigr)
\end{equation}
with linear maps
\begin{align}
  C_{L,\mathrm{eff}} &= C_L\bigl(1 + \tfrac{\partial C_L}{\partial z} z
    + \tfrac{\partial C_L}{\partial\mu}\mu\bigr),\\
  C_{D,\mathrm{eff}} &= C_D\bigl(1 + \tfrac{\partial C_D}{\partial z} z
    + \tfrac{\partial C_D}{\partial\mu}\mu\bigr),
\end{align}
clamped to $(0.2,3.0)$ so an NLP cannot invert aero. Defaults
$\partial C_L/\partial z = -8\,\si{\per\metre}$ (more downforce as the
car squats in the fastest-lap $z$ convention). Confirm the sign against
track testing before using a ride-height sweep as a design decision.

Front aero distribution
\[
  \eta_a = (x_{\mathrm{pc}}-x_r)/L \approx 0.48
\]
is a geometric lever, not a second $C_L$.

% ========================================================================
\section{Linear kinematics (not hardpoints)}
\label{sec:kin}
% ========================================================================

\begin{align}
  \gamma_{L,R} &= \gamma_0 \pm (\partial\gamma/\partial\phi)\,\phi,\\
  \delta_{\mathrm{toe},L,R} &= \delta_{t0} \pm (\partial\delta_t/\partial\phi)\,\phi,
\end{align}
plus front steer $\delta$ on the $Z$ rotation. Default $\gamma_0=-1^\circ$,
camber gain $0.3\,\mathrm{rad/rad}$. This is a first-order Taylor
expansion of a motion ratio, not a three-dimensional suspension.

% ========================================================================
\section{G-G diagrams}
\label{sec:gg}
% ========================================================================

Milliken's G-G diagram~\cite{milliken} is the set of $(a_x,a_y)$ the car
can sustain at speed $v$. Numerically, fastest-lap solves the 6DOF
steady equations on a grid of lateral accelerations and returns the
maximum and minimum longitudinal acceleration (throttle versus brake).

\paragraph{Speed dependence.}
Power-limited $a_x \sim P/(m v)$ falls with $v$. Downforce $\sim v^2$
raises $F_z$ and therefore $a_y^{\max}$. A single friction circle is
not an operating map; a stack of ellipses versus speed is.

\paragraph{Studio vehicle envelope} (\texttt{gg\_from\_vehicle}), used
when synthetic G-G is enabled (browser / no C++ library). $C_L$ here is
the fastest-lap (positive downforce) coefficient after the OpenVEHICLE
sign flip. Load factor $L$ is dimensionless:
\begin{align}
  L &= 1 + \frac{\tfrac12\rho \max(C_L,0) A v^2}{m g},\\
  a_y^{\max}/g &= \min(2.8,\; \mu_y L),\\
  a_{x,\mathrm{power}}/g &= \frac{P_{\max}/v - \tfrac12\rho |C_D| A v^2}{m g},\\
  a_{x,\mathrm{tire}}/g &= 0.65\,\mu_x L,\\
  a_x^{\max}/g &= \min(a_{x,\mathrm{power}}, a_{x,\mathrm{tire}}),\\
  a_x^{\min}/g &= -\min\bigl(\mu_x L,\; 2 T_{\mathrm{brake}}/(R m g)\bigr).
\end{align}
The factor $0.65$ is a driven-axle / combined-slip derate: full
$\mu_x F_z$ cannot be used for thrust while cornering. Two braked
axles at $T_{\mathrm{brake}}=\SI{1200}{\newton\metre}$ each give the
$2T/(Rmg)$ term ($R=\SI{0.203}{\metre}$). A friction ellipse
$a_x(a_y) = a_{x,\mathrm{peak}}\sqrt{1-(a_y/a_y^{\max})^2}$ is tabulated
at $v\in\{8,16,28,40\}\,\si{\metre\per\second}$. This envelope is
intentionally simpler than Pacejka G-G so that Studio mass and power
edits change lap time without \texttt{libfastestlapc}.

% ========================================================================
\section{Quasi-steady-state lap (OpenLAP three-pass)}
\label{sec:qss}
% ========================================================================

The racing line is the OpenTRACK polyline: curvature $\kappa(s)$ in
\si{\per\metre}. Point-mass tangent acceleration:
\begin{equation}
  a_y = v^2 \kappa \quad [\si{\metre\per\second\squared}],\qquad
  \frac{a_y}{g} = \frac{v^2\kappa}{g}.
  \label{eq:ay}
\end{equation}
Corner-limited speed: $v_{\max}=\sqrt{a_y^{\max}(v)/|\kappa|}$, iterated
because $a_y^{\max}$ depends on $v$.

Forward pass (drive envelope):
\begin{equation}
  v_{i+1}^2 = \min\bigl(v_{\max,i+1}^2,\; v_i^2 + 2 a_{x,\max} g\,\Delta s\bigr).
\end{equation}
Backward pass (brake envelope):
\begin{equation}
  v_i^2 = \min\bigl(v_{\max,i}^2,\; v_{i+1}^2 - 2 a_{x,\min} g\,\Delta s\bigr)
\end{equation}
with $a_{x,\min}<0$. Final
$v=\min(v_{\max},v_{\mathrm{acc}},v_{\mathrm{brk}})$. Time
\begin{equation}
  t_{i+1}=t_i+\frac{\Delta s}{\tfrac12(v_i+v_{i+1})}.
\end{equation}
Then
\begin{equation}
  a_x/g = \frac{v_{i+1}^2-v_i^2}{2\Delta s\, g}.
\end{equation}

\eng{Larger $|\kappa|$ lowers $v_{\max}$. A longer straight raises top
speed until power or drag limits. QSS cannot capture snap oversteer:
there is no yaw-inertia ODE on the track.}

\codepath{\texttt{qss\_lap.py} function \texttt{qss\_lap()}. Mesh from
\texttt{mesh\_opentrack}.}

% ========================================================================
\section{HUD channels: reconstructed, not integrated 6DOF}
\label{sec:hud}
% ========================================================================

After QSS, the dashboard requires throttle position (TPS), brake
position (BPS), steer, $\beta$, and four $F_z$. Those quantities are
not taken from \texttt{fsae6dof} states.

\paragraph{Pedals.}
\[
  \mathrm{TPS}=\mathrm{clip}(a_x/a_{x,\max},0,1),\quad
  \mathrm{BPS}=\mathrm{clip}(a_x/a_{x,\min},0,1)
\]
using the same envelope interpolation as the lap.

\paragraph{Steer and sideslip (OpenLAP bicycle).}
Kinematic Ackermann (degrees) is
$\delta_{\mathrm{kin}}=\arctan(L\kappa)$ converted with
\texttt{math.degrees}. Extra steer and body sideslip $\beta$ invert the
linear bicycle that produces centripetal force $m v^2\kappa$
(newtons; $\kappa$ in \si{\per\metre}):
\begin{gather}
  a=(1-d_f)L,\qquad b=-d_f L,\\
  \begin{pmatrix} 2C_F & 2(C_F+C_R)\\ 2 C_F a & 2(C_F a + C_R b)\end{pmatrix}
  \begin{pmatrix}\delta_{\mathrm{extra}}\\ \beta\end{pmatrix}
  =
  \begin{pmatrix} m v^2\kappa \\ 0 \end{pmatrix}.
\end{gather}
\texttt{bicycle\_steer} uses Cramer's rule:
$\delta_{\mathrm{extra}}=c_{11} b_0/\det$,
$\beta=-c_{10} b_0/\det$, $b_0=m v^2\kappa$. Then
$\delta=\delta_{\mathrm{extra}}+\delta_{\mathrm{kin}}$ (degrees) and
handwheel $=\delta\times$ rack ratio. $C_F,C_R$ must be
\si{\newton\per\radian} (convert if the MATLAB book is N/deg).

\paragraph{Lumped $F_z$ (newtons, downward).}
With $a_x,a_y$ already in $g$,
\begin{align}
  \Delta F_{\mathrm{long}} &= m (a_x g) h / L,\\
  \Delta F_{\mathrm{lat},f} &= m (a_y g) h d_f / t_f,\\
  \Delta F_{\mathrm{lat},r} &= m (a_y g) h (1-d_f) / t_r.
\end{align}
Front axle load $d_f m g - \Delta F_{\mathrm{long}}$ is split
$\pm\tfrac12\Delta F_{\mathrm{lat},f}$. A floor of \SI{50}{\newton}
avoids zeros in the HUD. The sign matches $m a_y h / t$; the magnitude
does not match the 6DOF spring network.

\paragraph{Power and energy on the tires card.}
$P \approx a_x g m v$ allocated by $F_z$ share; energy integrates
$|P|\,\mathrm{d}t$. These are HUD units, not inverter energy.

% ========================================================================
\section{Parameter provenance}
\label{sec:numbers}
% ========================================================================

\begin{table}[ht]
\centering
\caption{UBCO 2026 EV defaults (workbook $\to$ XML).}
\scriptsize
\begin{tabular}{@{}llll@{}}
\toprule
Quantity                 & Default        & Unit        & Provenance \\
\midrule
Total mass $m$           & 277.2          & kg          & VD OpenVEHICLE \texttt{Total Mass}; car+driver target. \\
Front mass dist.\ $d_f$  & 51.27          & \%          & VD book. \\
Wheelbase $L$            & 1620           & mm          & VD / chassis. XML splits $0.789$ and $-0.831$\,m. \\
Track $t_f,t_r$          & 1225, 1220     & mm          & VD FastestLap sheet. \\
CG height $h$            & 250            & mm          & VD estimate; XML \texttt{com.z=-0.25}. \\
$C_L,C_D$ (OV)           & $-3.77,-1.4$   & ---         & VD aero; flipped for fastest-lap. \\
$A$                      & 1.14           & m$^2$       & VD frontal area. \\
$\rho$                   & 1.2            & kg/m$^3$    & VD / near-ISA sea level. \\
$P_{\max}$               & 80             & kW          & FSAE EV rule~\cite{fsae-rules}. \\
$T_{\mathrm{peak}}$      & 240            & N$\cdot$m   & Powertrain / motor datasheet via VD. \\
Gear ratio                & 4.8            & ---         & Final drive in book. \\
Wheel rates               & 130k / 143.9k  & N/m         & Suspension / VD. \\
$\mu_x,\mu_y$            & 0.9, 1.5       & ---         & Interim; replace with TTC. \\
$I_{xx,yy,zz}$           & 30, 90, 110    & kg$\cdot$m$^2$ & Estimates in FastestLap comments. \\
2019 endurance map        & ---            & m           & VD OpenTRACK of FSAE 2019 endurance. \\
\bottomrule
\end{tabular}
\end{table}

The OpenVEHICLE book is the club source of mass, power, and geometry
~\cite{jones-linkedin,openlap}. Studio edits are experiments. Exporting
\texttt{.xlsx} (and saving to ranking) is how an experiment becomes a
candidate build file.

\paragraph{Mk~4 historical scale.}
The IC Mk~4 published \SI{245}{\kilo\gram}, wheelbase
\SI{1525}{\milli\metre}, $C_L A=\SI{2.6}{\metre\squared}$
~\cite{okmotorsports}. The 2026 EV book is a different vehicle; Mk~4
aero must not be copied onto EV-2 without aero sign-off.

% ========================================================================
\section{Code map}
\label{sec:code}
% ========================================================================

New software members can treat the following as the call graph of a
Studio calculate request.

\begin{lstlisting}[caption={Plant assembly (\texttt{fsae6dof.h}).}]
Tire_pacejka_std x 4
  -> Axle_car_6dof_fsae front  (STEERING_WITH_KAPPA)
  -> Axle_car_6dof_fsae rear   (POWERED_WITH_DIFFERENTIAL)
      -> Chassis_car_6dof_fsae (heave/roll/pitch, throttle, brake_bias,
                                pressure_center, aero maps, tire T)
          -> Road -> Dynamic_model_car
\end{lstlisting}

\begin{lstlisting}[language=Python,caption={Studio calculate (simplified).}]
params = derived_vehicle(read_openvehicle_xlsx(car_xlsx))
table  = gg_from_vehicle(params)   # or Driver envelope, or gg_diagram
mesh   = mesh_opentrack(map_xlsx)
result = qss_lap(mesh, table, v_cap=v_cap)
view   = reconstruct_lap(result, mesh, table, params=params)
write_hud_html(view, "hud.html")
write_results_html(view, "results.html")
\end{lstlisting}

Public type string: \texttt{fsae-6dof}. C++ class: \texttt{fsae6dof}.
Repository:~\cite{vehicle-model-v2}. Engine provenance:
~\cite{manzanero-fastestlap}.

% ========================================================================
\section{Limitations}
\label{sec:limits}
% ========================================================================

\begin{itemize}[leftmargin=1.6em]
  \item QSS HUD $F_z$ and steer are reconstructions. Do not size
        uprights from HUD $F_z$ alone.
  \item Pacejka shapes are placeholders until TTC~\cite{ttc}.
  \item QSS has no battery energy state; 80\,\si{\kilo\watt} is a power
        cap, not pack energy.
  \item No track elevation or banking.
  \item Vehicle-parameter G-G is an ellipse, not a 6DOF friction
        surface.
  \item Optimal lap time (transient NLP) is registered in fastest-lap
        but is not the Studio path.
\end{itemize}

% ========================================================================
\section{Onboarding summary}
\label{sec:onboard}
% ========================================================================

The sprung mass has six rigid-body degrees of freedom. Tires apply
contact forces; the series spring network shares vertical load
left/right; the motor is limited by peak torque or
80\,\si{\kilo\watt}, whichever is more restrictive at the present
speed; aerodynamic load grows with $v^2$ and with ride/pitch maps. A
Studio lap time is: follow the OpenTRACK line, respect the G-G
envelope in $a_y$ and $a_x$, integrate $\mathrm{d}s/v$. HUD steer and
sideslip are the OpenLAP bicycle reconstructed on that speed trace, so
existing MATLAB plots remain comparable~\cite{jones-linkedin,openlap}.

Change official numbers in the car \texttt{.xlsx}, run Calculate, then
compare lap time.

\bibliographystyle{plainnat}
\bibliography{ubco-2026-vehicle-model}

% ========================================================================
\appendix
\section{Selected symbols}
% ========================================================================

\begin{tabular}{@{}lll@{}}
\toprule
Symbol                    & Meaning              & SI unit \\
\midrule
$m$                       & Total mass           & kg \\
$L$                       & Wheelbase            & m \\
$t_f,t_r$                 & Track widths         & m \\
$h$                       & CG height            & m \\
$\kappa_{\mathrm{path}}$  & Path curvature       & 1/m \\
$\kappa_{\mathrm{tire}}$  & Longitudinal slip    & 1 \\
$\alpha$                  & Slip angle           & rad \\
$C_F,C_R$                 & Cornering stiffness  & N/rad \\
$P_{\max}$                & Motor power cap      & W \\
HUD                       & Heads-up display     & --- \\
QSS                       & Quasi-steady-state   & --- \\
\bottomrule
\end{tabular}

% ========================================================================
\section{File index}
\label{app:files}
% ========================================================================

\begin{description}[leftmargin=1.6em]
  \item[\texttt{src/core/vehicles/fsae6dof.h}]
        Vehicle type; G-G steady equations.
  \item[\texttt{src/core/chassis/chassis\_car\_6dof\_fsae.hpp}]
        Force/moment assembly, aero maps, thermal ODE.
  \item[\texttt{src/core/chassis/axle\_car\_6dof\_fsae.hpp}]
        $F_z$ network, $\kappa$, EV split, LSD, brakes.
  \item[\texttt{database/vehicles/fsae/ubco-2026-ev.xlsx}]
        VD parameter book.
  \item[\texttt{examples/python/fsae/qss\_lap.py}]
        Three-pass QSS.
  \item[\texttt{examples/python/fsae/qss\_channels.py}]
        HUD reconstruction.
  \item[\texttt{examples/python/fsae/qss\_job.py}]
        \texttt{gg\_from\_vehicle}.
  \item[\texttt{docs/superpowers/specs/2026-08-14-fsae-6dof-design.md}]
        Locked 6DOF scope.
\end{description}

\end{document}
